% Generated draft: requires mathematical and source review.
\documentclass[11pt]{article}
\usepackage[margin=25mm]{geometry}
\usepackage{fontspec}
\setmainfont{Latin Modern Roman}
\usepackage{amsmath,amssymb,mathtools,unicode-math}
\setmathfont{Latin Modern Math}
\usepackage{enumitem,needspace,xurl,hyperref}
\setlength{\parindent}{0pt}
\setlength{\parskip}{5pt}
\setlength{\emergencystretch}{3em}
\newcommand{\N}{\mathbb{N}}
\newcommand{\Z}{\mathbb{Z}}
\newcommand{\Q}{\mathbb{Q}}
\newcommand{\R}{\mathbb{R}}
\newcommand{\C}{\mathbb{C}}
\newcommand{\F}{\mathbb{F}}
\newcommand{\A}{\mathbb{A}}
\newcommand{\PP}{\mathbb{P}}
\newcommand{\E}{\mathbb{E}}
\newcommand{\eps}{\varepsilon}
\newcommand{\dd}{\,\mathrm{d}}
\begin{document}
\section*{UnsolvedMath ID: 30001680}
\textbf{Source problem code:} catalogue entry
\subsection{Definitions and mathematical statement}
Work in \(\mathsf{ZFC}\). For an uncountable regular cardinal \(\kappa\), use the following convention: an ideal on \(\kappa\) is a proper family \(I\subseteq\mathcal P(\kappa)\) containing every singleton, closed under subsets and unions of fewer than \(\kappa\) members. Thus it contains every bounded subset of \(\kappa\). The supplied extract does not explicitly specify this completeness convention; the questions below use it and require confirmation against the original report.

Write \(I^+=\mathcal P(\kappa)\setminus I\). Forcing with \(\mathcal P(\kappa)/I\) means forcing with the nonzero equivalence classes modulo \(I\), ordered by \([A]_I\leq[B]_I\) if \(A\setminus B\in I\). If \(G\) is generic over the ground universe \(V\), let
\[
U_G=\{A\in\mathcal P(\kappa)^V:[A]_I\in G\}.
\]
The ideal \(I\) is \emph{precipitous} if this forcing forces the generic ultrapower \(\operatorname{Ult}(V,U_G)\), formed using ground-model functions, to be well-founded. The ideal is \emph{normal} if, for every sequence \(\langle A_\alpha:\alpha<\kappa\rangle\) of members of \(I\), its diagonal union satisfies
\[
\mathop{\bigtriangledown}_{\alpha<\kappa}A_\alpha
=\{\beta<\kappa:\exists\alpha<\beta\ (\beta\in A_\alpha)\}\in I.
\]
A cardinal \(\delta\) is \emph{supercompact} if, for every cardinal \(\lambda\geq\delta\), there is an elementary embedding \(j:V\to M\) into a transitive class \(M\) such that \(\operatorname{crit}(j)=\delta\), \(j(\delta)>\lambda\), and \({}^{\lambda}M\subseteq M\), with closure computed in the ambient universe.

\textbf{Consistency questions.} Under these conventions, ask separately whether each of the following theories has a model:
\begin{enumerate}
\item \(\mathsf{ZFC}\) together with
\[
\begin{gathered}
\text{there exists a precipitous ideal on }\aleph_2,\\
\text{there exists no normal precipitous ideal on }\aleph_2.
\end{gathered}
\]
The second clause excludes every normal precipitous ideal on \(\aleph_2\), rather than merely requiring the exhibited ideal to be nonnormal.
\item \(\mathsf{ZFC}\) together with
\[
\begin{gathered}
\text{there exists a supercompact cardinal},\\
\text{there exists no precipitous ideal on }\aleph_1.
\end{gathered}
\]
\end{enumerate}
All cardinals, ideals, and quantifiers in each question are interpreted in the proposed model. The broader request to generalize a separation result beyond \(\aleph_1\) supplies motivation, but the extract does not state that antecedent result or a precise generalization for arbitrary higher cardinals. No particular relative-consistency hypothesis is supplied.
\subsection{Short English statement}
Ask whether a model of \(\mathsf{ZFC}\) can have a precipitous ideal on \(\aleph_2\) while having no normal precipitous ideal on that cardinal. Separately, ask whether a supercompact cardinal can coexist with the absence of precipitous ideals on \(\aleph_1\). The ideal conventions and the scope of the broader higher-cardinal generalization require source review.
\subsection{Sources}
\begin{itemize}
\item Local UnsolvedMath registry, supplied record with permanent ID \(30001680\), source problem code \texttt{catalogue entry}, title \emph{Precipitous Ideal Separation at Higher Cardinals}. Collection: Oberwolfach Reports Open Problems; proposed year: 2011. No individual proposer is identified in the supplied record.
\item Supplied report citation: \emph{Set Theory} (2011), Oberwolfach Reports, pp.~85--140, DOI \texttt{10.4171/owr/2011/02}. Source URL and literature-source label ``Oberwolfach problem statement'': \url{https://doi.org/10.4171/owr/2011/02}. The supplied record says the official report confirms the explicit higher-cardinal questions and assigns the statement status \texttt{corrected\_verified}; these are supplied verification claims, not independent verification here.
\item Supplied literature-source label: ``Gitik, precipitous ideals without normality'': \url{https://doi.org/10.1142/S0219061312500080}. The record supplies no fuller bibliographic details.
\item Literature assessment dated exactly \texttt{2026-08-21}: the supplied curation labels the problem open and states that the cited \(\aleph_1\) construction did not settle the higher-cardinal and supercompact variants, which it assesses as remaining open. This is supplied dated curation, not an independently verified assessment. The background explicitly requests independent verification of time-sensitive status and cited claims.
\end{itemize}
\end{document}
